\section{Validação}

\subsection{Estratégia de teste}

A validação foi realizada por meio de um \textit{harness} de testes automatizados integrado ao executável \texttt{genesys\_tools\_application}, acionado pelo comando \texttt{-{}-demo}. O harness instancia \texttt{DataAnalyzerDefaultImpl1} diretamente e executa 26 verificações organizadas em 7 seções, imprimindo \texttt{[PASS]} ou \texttt{[FAIL]} para cada uma.

Três conjuntos de dados fixos são usados:
\begin{itemize}
    \item \texttt{kNormalData}: 50 valores amostrados de $N(5, 1)$;
    \item \texttt{kExpoData}: 20 valores amostrados de $\text{Exp}(2)$ (média $= 2$);
    \item \texttt{kTemporalData}: 20 valores com tendência crescente, de 1,1 a 3,0.
\end{itemize}

\subsection{Seção 1 -- Carregamento de dados}

\begin{itemize}
    \item \texttt{setDataValues} aceita 50 valores e \texttt{summaryStatistics().n == 50};
    \item \texttt{setDataFilename} carrega os mesmos 50 valores de arquivo RawNumeric temporário;
    \item \texttt{clearData} esvazie o conjunto (\texttt{n == 0});
    \item \texttt{loadSecondSample} armazena um segundo vetor de 20 valores.
\end{itemize}

\subsection{Seção 2 -- Estatísticas descritivas}

\begin{itemize}
    \item Média dentro do intervalo $[4{,}5;\ 5{,}5]$ para dados $N(5,1)$;
    \item Ordenamento: $\text{mín} < Q_1 < \text{mediana} < Q_3 < \text{máx}$;
    \item Desvio-padrão $> 0$;
    \item Consistência entre \texttt{quartile(2)}, \texttt{decile(5)} e \texttt{centile(50)} (todos devem ser iguais à mediana).
\end{itemize}

\subsection{Seção 3 -- Histograma e boxplot}

\begin{itemize}
    \item Soma das frequências absolutas igual a $n = 50$;
    \item Histograma com exatamente 8 classes;
    \item Soma das frequências relativas igual a 1,0 (tolerância $10^{-10}$);
    \item $\text{IIQ} = Q_3 - Q_1$ (verificação da fórmula do boxplot).
\end{itemize}

\subsection{Seção 4 -- Ajuste de distribuições}

\begin{itemize}
    \item Ajuste normal válido (\texttt{FitResult.valid == true});
    \item Parâmetro de média do ajuste normal dentro de $[4{,}5;\ 5{,}5]$;
    \item \texttt{fitAll} retorna resultado válido;
    \item SSE do \texttt{fitAll} $\leq$ SSE do ajuste normal (é a melhor entre todas);
    \item Média do ajuste exponencial nos dados \texttt{kExpoData} dentro de $[1{,}0;\ 3{,}5]$ (esperada $\approx 2$).
\end{itemize}

\subsection{Seção 5 -- Testes de aderência}

\begin{itemize}
    \item Teste qui-quadrado não rejeita ajuste normal a $\alpha = 0{,}05$ (dados são normais);
    \item Teste KS não rejeita ajuste normal a $\alpha = 0{,}05$.
\end{itemize}

\subsection{Seção 6 -- Séries temporais}

\begin{itemize}
    \item \texttt{movingAverage(3)} produz $n - w + 1 = 18$ valores;
    \item Primeiro valor da média móvel igual à média dos três primeiros elementos ($(1{,}1 + 1{,}3 + 1{,}2)/3$);
    \item \texttt{autocorrelation(5)} retorna \texttt{maxLag + 1 = 6} valores;
    \item $\hat{r}(0) = 1{,}0$ por construção;
    \item $\hat{r}(1) > 0{,}5$ para dados com tendência crescente;
    \item \texttt{correlogram(5) == autocorrelation(5)} (implementação compartilhada).
\end{itemize}

\subsection{Seção 7 -- Inferência paramétrica}

\begin{itemize}
    \item IC 95\% para a média contém o verdadeiro valor 5,0;
    \item Teste t bilateral para $\mu_0 = 5{,}2$ não rejeita ($\mu_0$ próximo à média amostral de 5,206);
    \item Teste t bilateral para $\mu_0 = 10{,}0$ rejeita (valor muito afastado);
    \item Limite superior do IC de variância $>$ limite inferior;
    \item Teste de médias para duas populações (normal vs.\ exponencial) rejeita $H_0: \mu_1 = \mu_2$.
\end{itemize}

\subsection{Resultado da execução}

A saída do comando \texttt{-{}-demo} é apresentada abaixo (trecho):

\begin{verbatim}
============================================================
DataAnalyzer Demo — GenESyS tools
============================================================
...
1. DATA LOADING
------------------------------------------------------------
  [PASS] setDataValues accepted 50 values
  [PASS] setDataFilename loaded same 50 values from RawNumeric text file
  [PASS] clearData empties the dataset
  [PASS] loadSecondSample stores a second vector
...
============================================================
Demo complete.
============================================================
\end{verbatim}

Todos os 26 testes passam com a implementação atual.

\subsection{Problemas encontrados e corrigidos}

Durante o desenvolvimento foram encontrados e corrigidos três problemas:

\begin{enumerate}
    \item \textbf{Incompatibilidade de tipos no CDF de Erlang:} \texttt{std::llround} retorna \texttt{long long int}, enquanto a literal \texttt{1L} é \texttt{long int}. O compilador recusou \texttt{std::max(1L, std::llround(p2))} por tipos conflitantes. Corrigido com \texttt{std::max(1LL, std::llround(p2))}.

    \item \textbf{Métodos não-const de \texttt{ConfidenceInterval}:} Os métodos \texttt{inferiorLimit()} e \texttt{superiorLimit()} de \texttt{HypothesisTester\_if::ConfidenceInterval} não são declarados \texttt{const}, portanto variáveis locais \texttt{const auto ci} falham com erro de qualificador. Corrigido declarando as variáveis CI como \texttt{auto ci} (não-const).

    \item \textbf{Testes de duas populações usando sobrecargas de uma população:} As versões iniciais de \texttt{testProportionTwoSamples} e \texttt{testVarianceTwoSamples} omitiam o argumento \texttt{\_count2}, ativando as sobrecargas de uma população e tratando a estatística da segunda amostra como valor nulo. Corrigido passando \texttt{\_count2} como quarto argumento, selecionando as sobrecargas de duas populações corretas.
\end{enumerate}
